\section{Primary ideals}
\label{am-ch4-primary}
\begin{definition}[Primary ideal]
\label{am-def-primary}
\source{atiyah-macdonald-commutative-algebra:page-0059}\uses{}
An ideal $\mathfrak{q}\ne A$ is primary if $xy\in\mathfrak{q}$ and $x\notin\mathfrak{q}$ imply
$y^n\in\mathfrak{q}$ for some $n>0$.  Equivalently, every zero-divisor in $A/\mathfrak{q}$ is
nilpotent.  A primary ideal is $\mathfrak{p}$-primary when $\rad(\mathfrak{q})=\mathfrak{p}$.
\end{definition}
\begin{proposition}[Radical of a primary ideal]
\label{am-prop-4-1}
\dcref{4.1}
\source{atiyah-macdonald-commutative-algebra:page-0059}\uses{}
If $\mathfrak{q}$ is primary, $\rad(\mathfrak{q})$ is the smallest prime ideal containing $\mathfrak{q}$.
\end{proposition}
\begin{proof}
Put $\mathfrak{p}=\rad(\mathfrak{q})$.  If $(xy)^m\in\mathfrak{q}$, primaryness gives $x^m\in\mathfrak{q}$ or
$y^n\in\mathfrak{q}$ for some $n$, hence $x\in\mathfrak{p}$ or $y\in\mathfrak{p}$.  Thus $\mathfrak{p}$ is prime;
it contains $\mathfrak{q}$ and every prime containing $\mathfrak{q}$ contains all radicals.
\end{proof}
\begin{proposition}[Maximal radical implies primary]
\label{am-prop-4-2}
\dcref{4.2}
\source{atiyah-macdonald-commutative-algebra:page-0060}\uses{}
If $\rad(\mathfrak{a})$ is maximal, then $\mathfrak{a}$ is primary.  In particular every
power of a maximal ideal is primary.
\end{proposition}
\begin{proof}
The image of $\rad(\mathfrak{a})$ in $A/\mathfrak{a}$ is its nilradical.  By
Proposition~\ref{am-prop-1-8}, the quotient has only one prime ideal, so every
element is a unit or nilpotent; its zero-divisors are therefore nilpotent.
\uses{am-prop-1-8,am-def-primary}
\end{proof}
\begin{lemma}[Intersections of primary ideals]
\label{am-lem-4-3}
\dcref{4.3}
\source{atiyah-macdonald-commutative-algebra:page-0060}\uses{}
The intersection of finitely many $\mathfrak{p}$-primary ideals is $\mathfrak{p}$-primary.
\end{lemma}
\begin{proof}
The radical of an intersection is the intersection of the radicals, hence $\mathfrak{p}$;
if $xy$ lies in the intersection and $x$ is outside it, primaryness in one
component forces a power of $y$ into every component.
\end{proof}
\begin{lemma}[Colon ideals of a primary ideal]
\label{am-lem-4-4}
\dcref{4.4}
\source{atiyah-macdonald-commutative-algebra:page-0060}\uses{}
If $\mathfrak{q}$ is $\mathfrak{p}$-primary, then $(\mathfrak{q}:x)=(1)$ for $x\in\mathfrak{q}$; if $x\notin\mathfrak{q}$,
$(\mathfrak{q}:x)$ is $\mathfrak{p}$-primary; and if $x\notin\mathfrak{p}$, $(\mathfrak{q}:x)=\mathfrak{q}$.
\end{lemma}
\begin{proof}
The first and third assertions are immediate.  If $x\notin\mathfrak{q}$, then
$\mathfrak{q}\subseteq(\mathfrak{q}:x)\subseteq\mathfrak{p}$ and radicals are equal.  If $yz\in(\mathfrak{q}:x)$
with $y\notin\mathfrak{p}$, then $xyz\in\mathfrak{q}$ forces $xz\in\mathfrak{q}$, so
$z\in(\mathfrak{q}:x)$.
\end{proof}
\begin{definition}[Primary decomposition]
\label{am-def-primary-decomp}
\source{atiyah-macdonald-commutative-algebra:page-0060}\uses{}
A primary decomposition of $\mathfrak{a}$ is a finite intersection
$\mathfrak{a}=\bigcap_{i=1}^n\mathfrak{q}_i$ of primary ideals.  It is minimal when the
radicals $\mathfrak{p}_i=\rad(\mathfrak{q}_i)$ are distinct and no component can be omitted.
\end{definition}
\begin{theorem}[First uniqueness theorem]
\label{am-thm-4-5}
\dcref{4.5}
\source{atiyah-macdonald-commutative-algebra:page-0061}\uses{}
If $\mathfrak{a}=\bigcap_i\mathfrak{q}_i$ is a minimal primary decomposition with
$\mathfrak{p}_i=\rad(\mathfrak{q}_i)$, then the $\mathfrak{p}_i$ are exactly the prime ideals occurring as
$\rad(\mathfrak{a}:x)$ for $x\in A$.  Hence they depend only on $\mathfrak{a}$.
\end{theorem}
\begin{proof}
By Lemma~\ref{am-lem-4-4},
$(\mathfrak{a}:x)=\bigcap_i(\mathfrak{q}_i:x)$ and its radical is the intersection of those
$\mathfrak{p}_i$ for which $x\notin\mathfrak{q}_i$, with the other factors equal to $(1)$.
Prime avoidance (Proposition~\ref{am-prop-1-11}) shows every prime radical so
obtained is one of the $\mathfrak{p}_i$.  For each $i$, choose
$x_i\in\bigcap_{j\ne i}\mathfrak{q}_j\setminus\mathfrak{q}_i$ by minimality; then
$\rad(\mathfrak{a}:x_i)=\mathfrak{p}_i$.
\uses{am-lem-4-4,am-prop-1-11}
\end{proof}
\begin{definition}[Associated, isolated, and embedded primes]
\label{am-def-associated-primes}
\source{atiyah-macdonald-commutative-algebra:page-0061}\uses{}
The primes in Theorem~\ref{am-thm-4-5} are associated with $\mathfrak{a}$.  The
minimal members are isolated (or minimal) primes; the others are embedded.
\end{definition}
\begin{proposition}[Minimal primes over a decomposable ideal]
\label{am-prop-4-6}
\dcref{4.6}
\source{atiyah-macdonald-commutative-algebra:page-0061}\uses{}
Every prime containing a decomposable ideal $\mathfrak{a}$ contains a minimal prime
associated with $\mathfrak{a}$.  Thus the minimal associated primes are precisely
the minimal primes over $\mathfrak{a}$.
\end{proposition}
\begin{proof}
If $\mathfrak{p}\supseteq\mathfrak{a}=\bigcap_i\mathfrak{q}_i$, then
$\mathfrak{p}\supseteq\bigcap_i\mathfrak{p}_i$; Proposition~\ref{am-prop-1-11} gives
$\mathfrak{p}\supseteq\mathfrak{p}_i$ for some $i$, and descending among the finite set yields a
minimal one.
\uses{am-prop-1-11,am-def-primary-decomp}
\end{proof}
\begin{proposition}[Zero-divisors in a decomposable quotient]
\label{am-prop-4-7}
\dcref{4.7}
\source{atiyah-macdonald-commutative-algebra:page-0062}\uses{}
For a decomposable ideal $\mathfrak{a}$ with associated primes $\mathfrak{p}_i$,
\[
\bigcup_i\mathfrak{p}_i=\{x:(\mathfrak{a}:x)\ne\mathfrak{a}\}.
\]
In particular, if $(0)$ is decomposable, the zero-divisors are the union of
the associated primes and the nilradical is the intersection of the minimal
associated primes.
\end{proposition}
\begin{proof}
Apply Theorem~\ref{am-thm-4-5} to $A/\mathfrak{a}$ and use the characterization of
zero-divisors as the union of radicals of annihilators from
Proposition~\ref{am-prop-1-14}.  The minimal-prime intersection is the
nilradical by Proposition~\ref{am-prop-1-8}.
\uses{am-thm-4-5,am-prop-1-14,am-prop-1-8}
\end{proof}
\begin{proposition}[Localization of primary ideals]
\label{am-prop-4-8}
\dcref{4.8}
\source{atiyah-macdonald-commutative-algebra:page-0062}\uses{}
If $\mathfrak{q}$ is $\mathfrak{p}$-primary and $S$ is multiplicatively closed, then
$S^{-1}\mathfrak{q}=S^{-1}A$ when $S\cap\mathfrak{p}\ne\varnothing$.  If $S\cap\mathfrak{p}=\varnothing$,
$S^{-1}\mathfrak{q}$ is $S^{-1}\mathfrak{p}$-primary and contracts to $\mathfrak{q}$.  Thus localization
gives the corresponding bijection between primary and contracted primary
ideals.
\end{proposition}
\begin{proof}
If $s\in S\cap\mathfrak{p}$, some power of $s$ lies in $\mathfrak{q}$ and becomes a unit.  If
$S\cap\mathfrak{p}=\varnothing$ and $sa\in\mathfrak{q}$ with $s\in S$, primaryness gives
$a\in\mathfrak{q}$, so contraction is $\mathfrak{q}$; radicals localize by Proposition~\ref{am-prop-3-11},
and the defining zero-divisor condition is checked after clearing denominators.
\uses{am-prop-3-11}
\end{proof}
\begin{proposition}[Localized primary decompositions]
\label{am-prop-4-9}
\dcref{4.9}
\source{atiyah-macdonald-commutative-algebra:page-0063}\uses{}
Let $\mathfrak{a}=\bigcap_i\mathfrak{q}_i$ be minimal, with $\mathfrak{p}_i=\rad(\mathfrak{q}_i)$, and suppose
$S$ meets exactly $\mathfrak{p}_{m+1},\ldots,\mathfrak{p}_n$.  Then
\[
S^{-1}\mathfrak{a}=\bigcap_{i=1}^mS^{-1}\mathfrak{q}_i
\]
is a minimal primary decomposition, and contraction recovers $\mathfrak{a}$.
\end{proposition}
\begin{proof}
Localize the intersection using Proposition~\ref{am-prop-3-11}; the components
whose radicals meet $S$ become the unit ideal, and the others remain primary by
Proposition~\ref{am-prop-4-8}.  Distinct surviving radicals remain distinct.
\end{proof}
\begin{theorem}[Second uniqueness theorem]
\label{am-thm-4-10}
\dcref{4.10}
\source{atiyah-macdonald-commutative-algebra:page-0063}\uses{}
For an isolated set of associated primes, the intersection of the corresponding
primary components is independent of the chosen minimal primary decomposition.
\end{theorem}
\begin{proof}
Let $S=A\setminus\bigcup_{\mathfrak{p}\in\Sigma}\mathfrak{p}$.  Isolation and prime avoidance
ensure precisely the primes in $\Sigma$ miss $S$.  Proposition~\ref{am-prop-4-9}
shows the desired intersection is $S(\mathfrak{a})$, which depends only on
$\mathfrak{a}$ and $S$.
\uses{am-prop-4-9,am-prop-1-11}
\end{proof}
\begin{corollary}[Uniqueness of isolated components]
\label{am-cor-4-11}
\dcref{4.11}
\source{atiyah-macdonald-commutative-algebra:page-0063}\uses{}
The primary components corresponding to minimal (isolated) primes are uniquely
determined by $\mathfrak{a}$.
\uses{am-thm-4-10}
\end{corollary}
\begin{proof}
Apply Theorem~\ref{am-thm-4-10} to the set of minimal associated primes.
\end{proof}
